\chapter{Illustrative Examples}
\label{ch:examples}

The theorems of \cref{ch:theorems} become usable only after they have been seen
at work. Every computation below is written in full, in the order in which it
should be produced on an examination sheet.

\section{Rank computed in three ways}

\begin{example}{One matrix, three methods}{rank-three-ways}
Let
\[
  A=\begin{pmatrix} 1 & 2 & 3\\ 2 & 4 & 6\\ 1 & 1 & 1 \end{pmatrix}
  \in\Mat{3}{\R}.
\]

\textbf{By elimination.} Subtracting twice the first row from the second gives a
zero row; subtracting the first row from the third gives $(0,-1,-2)$. Hence
\[
  A\sim\begin{pmatrix} 1 & 2 & 3\\ 0 & -1 & -2\\ 0 & 0 & 0\end{pmatrix},
  \qquad \rank A=2 .
\]

\textbf{By the image.} The columns are $c_1=(1,2,1)$, $c_2=(2,4,1)$,
$c_3=(3,6,1)$. One checks $c_3=2c_2-c_1$, while $c_1,c_2$ are not proportional,
so $\im A=\Span(c_1,c_2)$ has dimension $2$.

\textbf{By the kernel.} Solving $Ax=0$: the second equation is twice the first,
and the system reduces to $x_1+2x_2+3x_3=0$, $x_2+2x_3=0$, whence
$x_2=-2x_3$ and $x_1=x_3$. Thus $\Ker A=\Span\bigl((1,-2,1)\bigr)$ has
dimension $1$, and $\rank A=3-1=2$ by \cref{thm:rank-nullity}.

The three methods must agree, and comparing them is the cheapest way to detect
an arithmetic slip.
\end{example}

\section{A projection, seen algebraically and geometrically}

\begin{example}{Decomposition induced by a projection}{projection-example}
Let $p\colon\R^{2}\To\R^{2}$ be given in the canonical basis by
\[
  P=\frac{1}{5}\begin{pmatrix} 4 & 2\\ 2 & 1\end{pmatrix}.
\]
Then $P^{2}=\frac{1}{25}\begin{pmatrix}20&10\\10&5\end{pmatrix}=P$, so $p$ is a
projection. Its image is $\Span\bigl((2,1)\bigr)$ and its kernel is
$\Span\bigl((1,-2)\bigr)$, since $P(1,-2)^{\mathsf{T}}=0$. These two lines are
orthogonal, so $p$ is the orthogonal projection onto the line $y=x/2$, exactly
the situation of \cref{fig:projection}.

By \cref{thm:split-simple} applied to $X^{2}-X=X(X-1)$, the map $p$ is
diagonalisable with $\spec(p)=\set{0,1}$; in the basis $\bigl((2,1),(1,-2)\bigr)$
its matrix is $\diag(1,0)$. Note also $\tr P=1=\rank P$, an instance of
\cref{prob:idempotent}.
\end{example}

\section{Diagonalisation without a characteristic polynomial}

\begin{example}{An involution}{involution}
Let $u\in\End(\R^{3})$ with matrix
\[
  M=\begin{pmatrix} 0 & 1 & 0\\ 1 & 0 & 0\\ 0 & 0 & -1\end{pmatrix}.
\]
A direct computation gives $M^{2}=I_3$, so $X^{2}-1=(X-1)(X+1)$ annihilates $u$.
Its roots are simple, hence $u$ is diagonalisable by \cref{thm:split-simple} and
\[
  \R^{3}=\Ker(u-\id)\oplus\Ker(u+\id)
        =\Span\bigl((1,1,0)\bigr)\oplus\Span\bigl((1,-1,0),(0,0,1)\bigr).
\]
No determinant was computed. This is the standard time-saving move in an
examination.
\end{example}

\section{When diagonalisation fails}

\begin{example}{A non-diagonalisable matrix}{jordan-block}
Let $N=\begin{pmatrix}0&1\\0&0\end{pmatrix}$. Its only eigenvalue is $0$, and
$\Ker N=\Span\bigl((1,0)\bigr)$ has dimension $1<2$, so $N$ is not
diagonalisable by \cref{thm:diagonalisable-criterion}. Observe that $X^{2}$
annihilates $N$: an annihilating polynomial that is split but has a
\emph{multiple} root gives no information. The hypothesis of simple roots in
\cref{thm:split-simple} cannot be dropped.
\end{example}

\section{The spectral theorem in practice}

\begin{example}{Orthogonal diagonalisation of a symmetric matrix}{symmetric-example}
Let $A=\begin{pmatrix}2&1\\1&2\end{pmatrix}$, which is symmetric. Then
$\det(A-\lambda I)=(2-\lambda)^{2}-1=(\lambda-1)(\lambda-3)$, so
$\spec(A)=\set{1,3}$ with eigenvectors $(1,-1)$ and $(1,1)$. These are
orthogonal, as \cref{thm:spectral} predicts. Normalising,
\[
  Q=\frac{1}{\sqrt2}\begin{pmatrix}1&1\\-1&1\end{pmatrix},\qquad
  \transpose{Q}AQ=\diag(1,3).
\]
Since both eigenvalues are positive, $A$ is positive definite by
\cref{cor:spd-eigen}, and
$\tr A\cdot\tr A^{-1}=4\cdot\bigl(1+\tfrac13\bigr)=\tfrac{16}{3}\geq4=n^{2}$,
illustrating \cref{prob:trace-inequality}.
\end{example}

\section{Visual explanation: what rank measures}

\begin{center}
\begin{tikzpicture}[scale=0.95,>=Stealth,font=\small]
  % source space
  \draw[dmzrule] (0,0) rectangle (3,2.4);
  \node[text=dmzsoft,font=\footnotesize] at (1.5,2.65) {$E$, $\dim E=3$};
  \draw[dmzgreen,fill=dmzgreen!12] (0.5,0.4) rectangle (1.4,1.9);
  \node[text=dmzgreen,font=\footnotesize] at (0.95,0.15) {$\Ker u$};
  % target space
  \draw[dmzrule] (5.5,0) rectangle (8.5,2.4);
  \node[text=dmzsoft,font=\footnotesize] at (7,2.65) {$F$};
  \draw[dmztheorem,fill=dmztheorem!12] (6.1,0.9) rectangle (7.9,1.6);
  \node[text=dmztheorem,font=\footnotesize] at (7,0.6) {$\im u$, $\rank u=2$};
  % arrow
  \draw[->,dmzaccent,line width=0.9pt] (3.2,1.2) -- (5.3,1.2)
        node[midway,above,text=dmzaccent,font=\footnotesize] {$u$};
  \draw[->,dmzsoft,dashed] (1.0,1.15) -- (6.9,1.25);
\end{tikzpicture}
\end{center}

The kernel is the information destroyed by $u$; the image is the information
kept. \Cref{thm:rank-nullity} says that the two amounts add up to $\dim E$:
nothing is created and nothing disappears.
